% ============================================================
%  Глава 5. Реализация и сравнение результатов
% ============================================================
\chapter{Реализация и сравнение результатов}
\label{ch:results}

В этой главе описаны проведённые эксперименты и сравнение четырёх
подходов к оценке rainbow-опционов:
\begin{enumerate}
  \item статическое хеджирование, см.~\cref{sec:res-static};
  \item динамическое хеджирование с ребалансировкой,
        см.~\cref{sec:res-dynamic};
  \item декомпозиция цены и обучение остатка,
        см.~\cref{sec:res-residual};
  \item прямое обучение модели цена $\to$ котировка,
        см.~\cref{sec:res-direct}.
\end{enumerate}
Сводные результаты приведены в~\cref{sec:res-summary}.

\section{Статическое хеджирование}
\label{sec:res-static}

Статический хедж строится в момент $t=0$ из ванильных опционов на
базовые активы, форварда и одного корреляционного свопа на горизонт $T$;
коэффициенты выбираются из условия совпадения первых производных
портфеля и rainbow-опциона по риск-факторам (см.~\eqref{eq:dyn-hedge}).
В дальнейшем состав портфеля не пересматривается.

\placeholderfig{12cm}{6cm}%
{Распределение остаточного P\&L статического хеджа на тестовом окне.}%
{fig:static-pnl}

\placeholdertab{Метрики ошибки статического хеджа (RMSE, MAE) по
страйкам и срокам $T$.}{tab:static-metrics}

\todoblock{Подставить итоговые цифры RMSE/MAE по статической модели.}

\section{Динамическое хеджирование с ребалансировкой}
\label{sec:res-dynamic}

Динамическая модель отличается от статической тем, что состав
хедж-портфеля пересматривается в каждый момент времени из дискретной
сетки $\{t_k\}_{k=0}^{M}\subset[0,T]$. На каждом шаге заново решается
система условий~\eqref{eq:dyn-hedge}, рассчитанная по обновлённым
ценам и волатильностям. При этом возникают \emph{транзакционные
издержки} на покупку/продажу активов: предполагается, что они
пропорциональны изменению нотачи позиции с коэффициентом $\tau$.

Зависимость остаточного $P\&L$ от частоты ребалансировки нелинейна:
слишком редкая ребалансировка увеличивает несовпадение греков, слишком
частая~--- транзакционные издержки.

\placeholderfig{12cm}{6cm}%
{Зависимость дисперсии остаточного P\&L от количества точек
ребалансировки $M$ при разных коэффициентах транзакционных издержек
$\tau$.}{fig:dyn-pnl-vs-M}

\placeholdertab{Метрики ошибки динамического хеджа по разным $M$ и
$\tau$.}{tab:dynamic-metrics}

\section{Обучение модели остатка $R$}
\label{sec:res-residual}

Для остаточной компоненты $R$ из~\eqref{eq:rainbow-decomp} обучается
нейронная сеть.

\todoblock{Зафиксировать архитектуру (например, MLP с 4-мя слоями по
128 нейронов, ReLU, нормализация входов), функцию потерь (Huber или
MSE) и оптимизатор. Указать размер обучающей и валидационной выборок.}

\placeholderfig{12cm}{6cm}%
{Кривые обучения: train/val loss по эпохам.}{fig:residual-train}

\placeholderfig{12cm}{6cm}%
{Распределение предсказанной и истинной остаточной компоненты $R$ на
тестовом множестве.}{fig:residual-scatter}

\placeholdertab{Метрики ошибки модели остатка $R$.}{tab:residual-metrics}

\section{Прямое обучение цены rainbow-опциона}
\label{sec:res-direct}

В качестве базовой линии (baseline) обучается модель того же класса, но
с целевой переменной~--- непосредственно цена rainbow-опциона. Это
позволяет ответить на вопрос: насколько выгодно <<выкидывать>> то, что
известно аналитически, прежде чем обучать нейронную сеть.

\placeholderfig{12cm}{6cm}%
{Распределение ошибок прямой модели против модели остатка.}%
{fig:direct-vs-residual}

\placeholdertab{Метрики ошибки прямой модели цены rainbow-опциона.}%
{tab:direct-metrics}

\section{Сводная таблица результатов}
\label{sec:res-summary}

\todoblock{Заменить плейсхолдеры цифрами после прогона полного
пайплайна.}

\begin{table}[H]
  \centering
  \caption{Сводные метрики ошибок всех методов на тестовой выборке.}
  \label{tab:summary-metrics}
  \begin{tabularx}{0.95\linewidth}{>{\raggedright\arraybackslash}X c c}
    \toprule
    \textbf{Метод} & \textbf{RMSE} & \textbf{MAE} \\
    \midrule
    Статическое хеджирование (\cref{sec:res-static})           & ---  & ---  \\
    Динамический хедж, M=? (\cref{sec:res-dynamic})            & ---  & ---  \\
    Модель остатка $R$ (\cref{sec:res-residual})                & ---  & ---  \\
    Прямая модель цены (\cref{sec:res-direct})                  & ---  & ---  \\
    \bottomrule
  \end{tabularx}
\end{table}

\todoblock{Прокомментировать: на каких страйках/корреляциях каждый метод
выигрывает, как растут ошибки при экстремальных $\rho\to\pm 1$, и какой
вклад в ошибку вносит ребалансировочный шум.}
